\begin{textsample}{POS Dim 1 – 2021 – Score 22.00 – 2021/t003892.txt}  \label{ex:f1_pos_062}
So \textbf{the} climate crisis, in my way of thinking about it, is \textbf{the} most serious manifestation of an underlying collision between human civilization as we ' ve presently organized it and \textbf{the} Earth ' s ecological systems.
And \textbf{the} system most in jeopardy is \textbf{the} very thin shell of atmosphere surrounding our \textbf{planet}, because we ' re spewing 162 million tons of human-caused global warming pollution into it every single day, as if this is an open sewer.
It ' s not an open sewer.
It ' s so thin that if you could drive an automobile at autobahn speeds to \textbf{the} top of that \textbf{blue} shell, you ' d reach it in about five minutes.
And that ' s where all \textbf{the} \textbf{greenhouse} gases are. \textbf{The} accumulated amount, coming from \textbf{the} burning of \textbf{fossil} fuels primarily, CO2 -- fastest-growing source of methane is from \textbf{fossil} fuels as well, also, \textbf{agriculture} as you have heard -- \textbf{the} accumulated amount traps as much heat now as would be released by 600, 000 Hiroshima-class atomic bombs exploding every day.
So \textbf{the} temperatures are going up at almost record levels every year.
Last year was \textbf{the} hottest year in recorded history, according to NASA, and \textbf{the} scientists say it ' s absolutely unequivocal that we are \textbf{the} cause of that.
And we ' re hearing Mother Nature and seeing \textbf{the} extreme events. \textbf{The} most anomalous extreme event since records began 200 years ago was in \textbf{the} Pacific Northwest of North America.
Forty nine and a half degrees, 121 degrees Fahrenheit in British Columbia.
In \textbf{the} tropics and subtropics, \textbf{the} combination of temperature and humidity is now making larger areas literally uninhabitable, \textbf{the} area so described is relatively small now, but if we continue over \textbf{the} next few decades, these areas are predicted to expand, and billions of people will be in areas where it is not safe to stay outside for more than a couple of hours.
Already, \textbf{the} climate refugees and migrants are four times more from \textbf{the} climate crisis than from all \textbf{the} wars and conflicts going on right now.
And \textbf{the} respected"Lancet"commission has predicted as many as a billion climate refugees in this century.
That ' s one of \textbf{the} reasons why we are seeing this wave of populist authoritarianism.
Sea level rise is also contributing, and of course, 93 percent of \textbf{the} extra heat is absorbed in \textbf{the} oceans, and \textbf{the} ocean temperatures are reaching record levels as well.
And that means more water vapor comes off \textbf{the} oceans, and \textbf{the} warmer ocean temperature makes \textbf{the} cyclonic storms, like hurricanes and typhoons much stronger.
Hurricane Ida struck \textbf{the} Gulf Coast as a category four, and as is so often \textbf{the} case, communities of color and poor people were disproportionately victimized.
This storm continued on north and northeast across North America to New York, New Jersey and Pennsylvania, killed a lot of people and dropped rain bombs.
And in New York City some basement apartments were flooded.
And I appreciate this family giving me \textbf{the} right to show this video.
He ' s safe, \textbf{the} guy is screaming for his mother.
He and his mother were both rescued by his brother, who punched a hole through \textbf{the} roof from \textbf{the} floor above.
These rain bombs are becoming much more frequent and much more extreme.
And actually \textbf{the} water coming off \textbf{the} oceans is held in much greater volumes by \textbf{the} warmer air.
Atmospheric rivers are becoming larger.
Here ' s one, from Hawaii in \textbf{the} lower left to Silicon Valley in \textbf{the} upper right, 2, 300 kilometers.
What was happening in Silicon Valley when this satellite picture was made?
Well, that ' s what happens when these rain bombs drop, and \textbf{the} average atmospheric river now contains as much moisture as 25 Mississippi Rivers.
They ' re creating what you could call"atmospheric \textbf{tsunamis}"in some areas. \textbf{The} downpours get much bigger and much more frequent, four times more frequent now than in 1980.
One of them dropped on Germany and neighboring countries in July, killed more than 200 people.
Look at \textbf{the} before and after.
This is an example of what we ' re doing to our \textbf{planet}.
Ten days ago, in part of Italy, there were 74 centimeters that fell in 12 hours, 29 \textbf{inches} of rain in 12 hours.
One year ago, October 3rd, as much rain fell here in \textbf{the} UK as there is water in Loch Ness.
And as a result, \textbf{the} insurance industry is now seeing record recoveries, and those recoveries I showed you from last year may get a tiny bit larger, depending on \textbf{the} outcome of this lawsuit. \textbf{The} owners and operators of this giant replica of Noah ' s Ark have sued their insurance company for a million dollars in flood damages.
Hard to make some of this stuff up. \textbf{The} same extra heat is pulling \textbf{the} moisture out of \textbf{the} first meter of \textbf{the} soil, creating \textbf{the} worst \textbf{droughts} in memory. 93 percent of \textbf{the} American West is in \textbf{drought} now, 100 percent of California, half of California is in \textbf{the} most extreme form of \textbf{drought}, and \textbf{the} same heat is also draining \textbf{the} reservoirs, evaporating \textbf{the} reservoirs.
Lake Mead is down to levels not seen since they started filling it in \textbf{the} 1930s.
In Brazil, \textbf{the} same thing is happening.
In Eastern Europe, \textbf{the} Czech Republic has been going through \textbf{the} worst \textbf{drought} in at least 500 years.
And in \textbf{the} southern cone of Africa, 100 million people are experiencing food insecurity primarily because of \textbf{the} extended \textbf{droughts}.
And when \textbf{the} temperatures go up, \textbf{the} fires get a lot worse and \textbf{the} worst fires in California and western US history have been in \textbf{the} last two years.
Also in Siberia, in Australia, in southern Europe.
It doesn ' t have to ruin your \textbf{golf} game.
And this is a serious point, because we can ' t let these conditions become \textbf{the} new normal.
It ' s not fine.
And this is an example of why it ' s not fine.
A lightning strike hitting a natural gas leak in \textbf{the} middle of \textbf{the} Gulf of Mexico.
Lightning gets more common with \textbf{the} climate crisis.
We ' re relying on dead plants and animals, leaving their residue in \textbf{the} atmosphere to threaten \textbf{the} species that are still alive, and 50 percent of them are in danger of extinction in this century.
And as we push into previously wild areas of \textbf{the} world, we encounter millions of new viruses that we have not dealt with in \textbf{the} past.
Five new infectious diseases every year, emerging, three quarters of them from animals, like \textbf{the} COVID-19 pandemic.
Which raises some of \textbf{the} same questions that are raised by \textbf{the} climate crisis.
When \textbf{the} world ' s leading scientists are setting their hair on fire to get our attention, should we listen to them?
Check.
Can our interconnected global civilization suddenly be turned upside down?
Check.
Are \textbf{the} poor and marginalized populations of \textbf{the} world \textbf{the} most affected?
Check.
Can science and technology give us nearly miraculous solutions in record time?
Check.
Will we deploy those solutions in time?
That is \textbf{the} question.
We have inequitable vaccine distribution in \textbf{the} world, and it threatens everyone in \textbf{the} world.
We also have solutions to \textbf{the} climate crisis, but they ' re not equitably distributed.
Worldwide, 90 percent of all of \textbf{the} new electricity generation installed last year was renewables.
Almost all of it from solar and wind.
Just one year before \textbf{the} Paris Agreement, solar and wind electricity was cheaper than \textbf{fossil} electricity in only one percent of \textbf{the} world.
Five years later, two thirds of \textbf{the} world.
Three years from now, it will be \textbf{the} cheapest source in 100 percent of \textbf{the} world. \textbf{Coal} is not getting cheaper.
Gas is not getting cheaper.
Nuclear has been getting more expensive.
Wind is cheaper than all three, and solar is continuing to plummet in cost and is \textbf{the} cheapest of all now because \textbf{the} cost of making \textbf{the} solar panels and \textbf{the} windmills continues to come down.
By \textbf{the} way, \textbf{the} \textbf{famous} \textbf{coal} museum in Kentucky just installed solar panels on its roof in order to save in its operating budget.
Now we ' re getting battery storage.
I ' ll rescale this graph to show you \textbf{the} predicted emergence of a new one-trillion-dollar industry in \textbf{the} next few decades and green hydrogen on top of that. \textbf{The} car batteries are getting cheaper as demand continues to increase.
And within less than two years in some model categories, \textbf{the} EVs will be cheaper than their internal combustion engine counterparts and within four years in all model categories.
So, \textbf{the} \textbf{oil} and gas industry has been \textbf{the} worst investment in global markets for more than a decade, while \textbf{the} clean energy companies are really becoming more profitable.
Now \textbf{the} \textbf{oil} and gas companies say that they ' re going to invest a lot more in renewables and CCS.
Yes, it ' s \textbf{the} fastest-growing -- they ' ve tripled their investments all \textbf{the} way up to four percent.
Ninety-six percent of \textbf{the} money they ' re spending is still going into \textbf{oil} and gas and \textbf{coal}. \textbf{The} impression they ' ve given us is not an honest impression, not for \textbf{the} first time.
And they ' re telling a different story to Wall Street.
They still have planes, trains, a few automobiles and ships, but they ' re telling Wall Street and financial markets they ' re going to make it up in plastics, which is 75 percent of their third-largest market, petrochemicals.
How ' s that working out for us?
Not so well.
Banning single-use plastics is one of \textbf{the} many things that we have to do.
We have to shift to \textbf{sustainability}.
And here ' s \textbf{the} good news, \textbf{the} \textbf{sustainability} revolution is \textbf{the} single biggest business opportunity in \textbf{the} history of \textbf{the} world.
It has \textbf{the} scale of \textbf{the} industrial revolution coupled with \textbf{the} speed of \textbf{the} digital revolution.
And what we ' re seeing is that we have to seize this opportunity, but we need reforms in \textbf{the} current version of capitalism.
Capitalism does a lot of things great.
It balances supply and demand and allocates resources efficiently and often \textbf{unlocks} a higher \textbf{fraction} of \textbf{the} human potential.
And by \textbf{the} way, \textbf{the} alternatives on \textbf{the} left and right explored in \textbf{the} 20th century didn ' t work out very well.
But in order to see how we need to change capitalism, we have to pull out our focus.
And let me present an analogy. \textbf{The} electromagnetic spectrum from \textbf{the} long radio waves to \textbf{the} short gamma rays and all \textbf{the} things in between, \textbf{the} portion of visible light that we can see with our eyes makes up only 0. 1 percent of \textbf{the} total.
For \textbf{the} eight years I worked in \textbf{the} White House, I got a daily report from \textbf{the} intelligence community that collected from all parts of \textbf{the} spectrum, and it was a much more accurate picture.
Here ' s \textbf{the} analogy. \textbf{The} value spectrum is something we too frequently look at through \textbf{the} very narrow aperture of short-term profits for one stakeholder, \textbf{the} shareholders.
But this leaves out negative externalities, as \textbf{the} economists call them, like \textbf{the} pollution.
That ' s why we need a price on \textbf{carbon} and a price on plastic pollution.
It also leaves out positive externalities.
So we ' re chronically underinvesting in education and health care and environmental protection, pandemic preparedness.
We ' re ignoring \textbf{the} depletion of resources like topsoil and \textbf{underground} water aquifers.
And we ' re ignoring \textbf{the} distribution of incomes and net worths to \textbf{the} point where one percent of \textbf{the} people have almost half of \textbf{the} world ' s wealth.
That ' s one of \textbf{the} other factors that ' s driving populist authoritarianism.
We have to take account of \textbf{the} environmental effects and \textbf{the} effects on people and their families and \textbf{the} communities where they live and \textbf{the} communities in their supply chain and \textbf{the} ethics in \textbf{the} C-suites.
And we have to realize that hyper inequality is a threat to both capitalism and to democracy.
There is an emergent form called multistakeholder capitalism, and it is driving a lot of new decisions.
For example, in asset management, almost half of all \textbf{the} assets in \textbf{the} world under management are now in portfolios committed to net-zero.
One reason is that \textbf{the} Paris Agreement set \textbf{the} direction of travel, where every country in \textbf{the} world committed to net-zero.
So just last May, \textbf{the} G7 nations banned financing of \textbf{coal} plants overseas.
And just three weeks ago, China did, which is good, because they were financing most of them.
I hope that they will cut down on \textbf{the} \textbf{coal} chain in China domestically as well.
And now more than half of all \textbf{the} \textbf{greenhouse} gas \textbf{emissions} and two thirds of global GDP is coming from countries that have set net-zero targets.
So here ' s \textbf{the} hope.
Once we reach net-zero, then with a lag time of as little as three to five years, \textbf{the} temperatures in \textbf{the} world will stop going up.
And once we reach net-zero within as little as 25 to 30 years, half of \textbf{the} human-caused CO2 will fall out of \textbf{the} atmosphere.
It is as if we have a switch that we can flip in order to stop \textbf{the} climate crisis.
Regrettably, some damage has been done, but we can stop \textbf{the} temperatures from going up and start \textbf{the} healing process.
But we all have to flip this switch known as reaching net-zero. \textbf{The} young people are telling us that we have to, and they ' re marching in every country in \textbf{the} world.
And what ' s \textbf{the} response?
As recently as this morning we heard,"Well, it may be impossible."Well, it ' s not impossible.
Nelson Mandela said,"It always seems impossible until it ' s done."That ' s what they told \textbf{the} abolitionists."It ' s impossible to eliminate slavery."That ' s what they told \textbf{the} women ' s suffragettes."It ' s impossible to have equal rights for women."In \textbf{the} civil rights movement in my country and \textbf{the} anti-apartheid movement in South Africa and more recently, \textbf{the} lesbian and gay liberation and equal rights movement.
We can do this.
This is \textbf{the} biggest emergent social movement in all of history.
We can do this.
And if anybody thinks that we don ' t have \textbf{the} political will, remember, political will is itself a \textbf{renewable} resource.
Thank you very much.
Thank you.

% matched lemmas: agriculture, blue, carbon, coal, drought, emission, famous, fossil, fraction, golf, greenhouse, inch, oil, planet, renewable, sustainability, the, tsunami, underground, unlock
\end{textsample}
